\section{Introduction}
\label{sec:intro}

The modern deep learning paradigm increasingly relies on pre-training large foundation models and subsequently fine-tuning them on specific downstream tasks~\cite{radford2021learning, devlin2018bert}. 
While this approach yields outstanding specialized experts, deploying a fleet of individual models for every unique task is computationally prohibitive in multi-task scenarios. 
To alleviate this burden, \emph{model merging} has emerged as an attractive, zero-shot framework to fuse multiple specialized weights into a single multi-task model without any joint retraining~\cite{ilharco2022editing, wortsman2022model, matena2022merging}.

However, despite rapid advancements, the model merging literature typically operates within a standard paradigm: Euclidean linearity. 
Virtually all existing techniques, ranging from simple task arithmetic~\cite{ilharco2022editing} and weighted Fisher merging~\cite{matena2022merging}, to advanced test-time adaptors like SyMerge~\cite{jung2025symerge}, operate under the fundamental assumption that task parameter spaces exist in flat, linearly-interpolated representations. 
Even manifold-based formulations such as OrthoMerge~\cite{yang2026orthomerge} apply predefined, rigid projections. 
These methods implicitly assume that the optimal path connecting disparate expert parameter basins is a straight line, or a projection thereof.

\begin{figure*}[t]
  \centering
  \includegraphics[width=0.95\textwidth]{symerge-plot.png}
  \caption{Conceptual overview of the FoldMerge (Neural Origami) paradigm. (Left) Traditional Euclidean model merging (e.g., Task Arithmetic, SyMerge) linearly averages task parameter vectors $\theta_1, \theta_2$, forcing the merged model through high-loss Euclidean barriers. (Right) FoldMerge learns a non-linear weight-space diffeomorphism $g_\phi$ that bends and deforms the underlying coordinate system, folding separate low-loss basins together into a single unified basin in Origami Space ($z$-space). The latent additive combination $\bar{z}$ is computed in this folded space and decoded back via $g_\phi^{-1}$, finding a non-linear low-loss path.}
  \label{fig:origami_concept}
\end{figure*}

This assumption is worth examining in light of our understanding of neural loss landscapes. 
Deep weight spaces are highly non-convex, curved, and complex~\cite{garipov2018loss, draxler2018essentially}. 
Downstream task-specific models lie in distinct, disjoint basins of attraction. 
As illustrated in Figure~\ref{fig:origami_concept}, forcing a straight-line Euclidean interpolation between these disjoint basins can drag the merged model across high-loss barriers. 
This can manifest empirically as task interference and accuracy degradation, particularly on fine-grained and heterogeneous tasks.

In this work, we explore an alternative to this flat-space assumption. 
We investigate the following fundamental question: \textbf{Can we model-merge by warping and transforming the weight-space coordinate system itself?}

We introduce \textbf{FoldMerge (Neural Origami)}, an exploratory, non-linear coordinate-transformation framework for model merging. 
Instead of searching for paths of low loss in a rigid, fixed Euclidean coordinate system, we train a differentiable coordinate-warping diffeomorphism $g_\phi: \mathbb{R}^d \to \mathbb{R}^d$ parameterized by $\phi$ using normalizing flows (specifically RealNVP~\cite{dinh2016density}). 
FoldMerge investigates the use of a continuous coordinate warp to map the pre-trained weights $\theta_{base}$ and task expert weights $\{ \theta_k \}_{k=1}^K$ into a latent, unified ``Origami Space'' where we perform our merging: $z_{base} = g_\phi(\theta_{base})$ and $z_k = g_\phi(\theta_k)$. 
We then compute the additive latent coordinates in Origami Space, $\bar{z} = z_{base} + \sum_k \lambda_k z_k$, and decode the merged model via the analytical inverse diffeomorphism: $\theta_{MTL} = g_\phi^{-1}(\bar{z})$. 
Because the coordinate system itself is warped, our straight-line combination in Origami Space maps back to a non-linear path in the original weight-space, aiming to mitigate Euclidean ridges of destructive interference.

To ensure that the coordinate warp remains smooth and stable, we introduce an implicit flow regularization penalty using parameter-wise $\ell_2$ weight decay on the normalizing flow network. 
This forces the diffeomorphism $g_\phi$ to be locally volume-preserving and to deviate minimally from an identity mapping, except where necessary to align the basins. 
Crucially, this implicit regularizer completely avoids any prohibitive $O(d^2)$ computational overhead of evaluating high-dimensional Jacobian matrices, making test-time adaptation computationally feasible. 
Following the self-labeling principles of SyMerge~\cite{jung2025symerge}, we optimize the flow parameters $\phi$ and merging coefficients $\lambda$ at test-time using unlabeled data streams guided by expert teacher predictions.

To investigate our hypothesis, we apply FoldMerge to the visual projection weight matrix (`model.visual.proj`) of a ViT-B/32 backbone on the standard 8-task classification benchmark. 
Evaluated under standard test-time adaptation settings, FoldMerge achieves an average accuracy of \textbf{89.76\%}, performing on par with the highly optimized SyMerge baseline~\cite{jung2025symerge} (89.74\%) while outperforming it on 5 out of 8 individual tasks. 
This provides an encouraging proof-of-concept for the viability and trainability of non-linear parameter warping.

In summary, our key contributions are:
\begin{itemize}
    \item \textbf{A New Geometric Perspective:} We propose FoldMerge, an exploratory framework to formulate model merging as a non-linear parameter-space warping process via learned weight-space diffeomorphisms, departing from the rigid Euclidean averaging paradigm.
    \item \textbf{Differentiable Coordinate Warping:} We implement a sequence of RealNVP affine coupling layers with bounded scale functions and implicit weight-decay regularization to investigate the mathematical viability and optimization behavior of learned coordinate warping.
    \item \textbf{Empirical Proof-of-Concept:} We demonstrate that non-linear parameter-space warping is computationally trainable, achieving an average accuracy of \textbf{89.76\%} on the 8-task ViT-B/32 benchmark and performing on par with highly optimized linear adaptive baselines.
    \item \textbf{Analysis of Structural Limitations:} We provide a transparent discussion of the practical and mathematical limitations of this approach—including coordinate-dependence, slicing category errors, and parameter/computational overhead—to serve as a foundation and guide for future research in non-linear weight alignment.
\end{itemize}
